\begin{titelseite}
\quellseite{2}
Dr. Martin Luthers\\[1em]
\werktitel{Sämmtliche Schriften,}\\[1em]
\gesperrt{herausgegeben von}\\[0.5em]
Dr. Joh. Georg Walch.\\[1em]
\werkuntertitel{Dritter Band.}\\[0.5em]
Auslegung des Alten Testaments.\\
(Fortsetzung.)\\[1em]
Neue revidirte Stereotypausgabe.\\[1em]
St. Louis, Mo.\\
CONCORDIA PUBLISHING HOUSE.\\
1894.
\end{titelseite}

\begin{titelseite}
\quellseite{3}
Dr. Martin Luthers\\[1em]
\werktitel{Auslegung des Alten Testaments.}\\[0.5em]
(Fortsetzung.)\\[1em]
Predigten über das erste Buch Mosis und Auslegungen über die folgenden
biblischen Bücher bis zu den Psalmen (excl.).\\[1em]
Aufs Neue herausgegeben im Auftrag des Ministeriums der deutschen
ev.-luth. Synode von Missouri, Ohio und anderen Staaten.\\[1em]
St. Louis, Mo.\\
CONCORDIA PUBLISHING HOUSE.\\
1894.
\end{titelseite}

\quellseite{5}
\chapter{Vorwort}

Nachdem in den beiden ersten Bänden unserer Ausgabe „Luthers große Auslegung des ersten Buchs Mosis“ mitgetheilt worden ist, folgen im gegenwärtigen Bande „Luthers Predigten über das erste Buch Mosis“ aus den Jahren 1523 bis 1526, sowie die Auslegungen der folgenden Bücher des Alten Testaments bis zu den Psalmen exclusive. Der Text aller dieser Schriften, welche im Inhaltsverzeichnisse angegeben sind, ist nach Kräften verbessert und berichtigt worden, und in der ersten Anmerkung zu jeder Schrift die nöthige Nachricht gegeben über die Veranlassung, Zeit der Abfassung, Veröffentlichung derselben etc. Einer besonderen Einleitung bedarf es daher nicht.

Durch die Lutherfunde in der Zwickauer Rathsschulbibliothek wurden wir in den Stand gesetzt, für einige der Schriften genauere und richtigere Zeitbestimmungen zu treffen, als dies in den bisherigen Ausgaben geschehen konnte, z.\,B. für Luthers Predigten über das erste Buch Mosis, für die Predigten über 1 Mos. 3, 15. und über 1 Mos. 22, 18.; für die Anmerkungen zu dem fünften Buch Mosis; die Zeitangaben für andere Schriften, z.\,B. „Luthers Auslegung über etliche Capitel des zweiten Buchs Mosis“, erhalten durch ebendenselben Fund ihre Bestätigung.

Zwei größere Schriften dieses Bandes haben wir neu aus dem Lateinischen übersetzt, nämlich (nach der Weimarischen Ausgabe): „Die zehn Gebote dem Volk zu Wittenberg gepredigt“ und (nach der lateinischen Erlanger Ausgabe) „Luthers Anmerkungen zu dem fünften Buche Mosis“. Wiewohl in der ersteren Schrift noch nicht das volle, reine Licht des Evangelii in allen Stücken leuchtet, so ist sie doch auch reich an heilsamer Lehre und äußerst interessant durch die ausführliche Darlegung der Sitten und Gebräuche im Zeitalter der Reformation, und hat großen Werth als reformationsgeschichtliches Actenstück. Die alte Uebersetzung der letzteren Schrift ist voll von sinnstörenden Fehlern, die theils auf Rechnung des Uebersetzers, theils auch wohl des Druckers zu setzen sind; z.\,B. (nach der alten Ausgabe) Col. 2294, § 13 „fleischlich“ statt: fälschlich (perverse
\quellseite{6}
); Col. 2323, § 30: Die Eltern aber, die ihr [der Jungfrau] kräftiglich widerstehen, statt: die ihm [Christo] starken Widerstand leisten; Col. 2332, § 16: es lautet grausamlich, statt: es deutet auf Grausamkeit hin (atrociter sonat); Col. 2336, § 25: daß die Stiefmutter den Sohn nicht nehmen soll, statt: daß man des Vaters Weib nicht nehmen soll (uxorem patris non accipere); Col. 2347, § 8: das bald folgende Gesetz, statt: das eben vorhergehende Gesetz; Col. 2349, § 13: der eine neue Pflanze des Glaubens ist, statt: ein Neuling im Glauben. Col. 2349 f., § 13: die Kinder des Bräutigams, statt: die Hochzeitsleute (filii sponsi); Col. 2405, § 44: „meine Feinde“ statt: Pfeile; Col. 2412, § 24: „vom Golde“ statt: vom Solde (mercede); Col. 2418, § 4: verführet durch falsche Aehnlichkeit der Wörterlein, statt: getäuscht durch die Aehnlichkeit des Ausdrucks, u.\,dgl.\,m.

Wiewohl „Luthers Vorlesung über das Buch der Richter“\footnote{„Dr. Martin Luthers Vorlesung über das Buch der Richter. Aus einer in der Zwickauer Rathsschulbibliothek befindlichen Handschrift herausgegeben von Georg Buchwald, Dr. Phil. Cand. Theol., Oberlehrer am Gymnasium zu Zwickau. Leipzig, Julius Dreschers Verlag. 1884.“ Ehe diese Schrift in die Weimarische Ausgabe aufgenommen wurde, hat Kawerau den Text noch einmal mit der Handschrift verglichen und verbessert. Buchwald selbst hatte sie l.\,c. S. 10 in das Jahr 1529 oder 1530 gesetzt, dagegen Köstlin in seiner Vorrede zu Buchwalds Edition läßt die Vorlesungen im Jahre 1516 oder 1517 gehalten sein. Letzterer Annahme schließt sich die Weimarische Ausgabe an, sagt jedoch (Nachträge zu Bd. IV), daß auch der von ihr angenommene Anfangstermin (1516) unsicher sei. So viel steht fest, es ist eine vorreformatorische Schrift. Ueber die Handschrift sagt die Weimarische Ausgabe: „Es ist eine saubere, aber mit wenig Verständniß nach einem Collegienhefte angefertigte Reinschrift. Die Handschrift erstreckt sich nicht gleichmäßig über sämmtliche Capitel des Richterbuches; sie bricht in der Auslegung von Cap. 18 ab [Cap. 15 und 17 fehlen ganz], zeigt auch sonst Lücken und trägt einen fragmentarischen Character; es war wohl eine unvollständige und ungleich gefertigte Nachschrift, die dem Schreiber zur Verfügung stand. Die zahlreichen Fehler im Manuscripte legen die Vermuthung nahe, daß der Schreiber nicht selber der Hörer der Vorlesung gewesen war, sondern sich nur auf die fehlerhafte Nachschrift eines Anderen für seine Arbeit angewiesen sah.“ An einer anderen Stelle, unter den „Nachträgen“ zum 4. Bande der Weimarischen Ausgabe zu S. 569, Anm. 1 (so ist zu lesen statt: „Zu S. 610, Anm. 1.“), wird über ein Citat aus einem Briefe des Erasmus an Nicol. Everardus, 15. Mai 1524, bemerkt: „Die Uebereinstimmung ist so groß, daß sie nur auf Entlehnung beruhen kann, und diese kann … nur auf Rechnung des Abschreibers der Richtervorlesung gesetzt werden, da die Vorlesung selbst unbedingt weit früher gehalten ist, als dieser Brief des Erasmus geschrieben resp. veröffentlicht wurde (1524 resp. 1529). Ist aber hier die Annahme geboten, daß der Verfasser der Zwickauer Abschrift … an dieser einen Stelle Lücken in dem ihm vorliegenden Collegienheft auf eigene Hand ergänzt hat, so wird höchst wahrscheinlich, daß auch die mehrfachen aus Augustin entlehnten resp. verwirrten Stücke auf Rechnung dieses Schreibers zu setzen sind, ja vielleicht auch der so stark an Luthers Decem praecepta anklingende Passus auf S. 538.“ — Wir sind fest überzeugt, daß der eben erwähnte Passus aus „Die zehn Gebote dem Volk zu Wittenberg gepredigt“ entlehnt ist. Vgl. Col. 1218 dieses Bandes.} Aufnahme gefunden hat in der Weimarischen Ausgabe, Bd. IV, S. 527, haben wir uns doch nicht entschließen können, dieser Schrift einen Platz einzuräumen unter den „Auslegungen Luthers über das Alte Testament“, wegen ihrer fragmentarischen und auch sonst zweifelhaften Beschaffenheit. Sollte aber das ehrwürdige Ministerium der Synode von Missouri, Ohio und anderen Staaten es wünschen, so werden wir sie einem Supplementbande einfügen.

Bei der letzten Schrift dieses Bandes, „Von den letzten Worten Davids“, einer
\quellseite{7}
der bedeutendsten Schriften Luthers, haben mir durch die Vergleichung der lateinischen Uebersetzung Crucigers einige Textesstellen berichtigen, andere dem Verständniß näher bringen können.

GOtt gebe in Gnaden auch zu dieser Arbeit seinen Segen, daß wir durch dieselbe unseren lieben Heiland, den HErrn JEsum Christum, immer besser erkennen und, wie Luther sagt (Col. 1959, § 139), „also den Vater und den Heiligen Geist, Einen GOtt, ewiglich sehen mögen. Wer den Sohn hat, dem steht die Schrift offen, und je größer und größer sein Glaube an Christum wird, je heller die Schrift ihm scheint“. Amen.

\begin{unterschrift}
St. Louis, um Michaelis 1894.\\
A. F. Hoppe.
\end{unterschrift}
